\chapter{Onglet Résultats}
\label{chap:resultats}

L'onglet \menu{Résultats} présente les résultats de simulation
sous forme d'une \textbf{grille configurable de cellules
d'analyse}. Chaque cellule affiche un type de tracé indépendant
(polaire, distribution de pression, profil de couche limite,
etc.).

\screenshot{10_tab_results_empty.png}{1.0}{%
    Onglet \menu{Résultats} avant toute simulation. La grille
    est par défaut en disposition $2\times2$, soit
    quatre cellules vides.}

\screenshot{18_tab_results.png}{1.0}{%
    Onglet \menu{Résultats} après simulation du profil courant
    (bleu), de la référence (rouge) et du profil avec volet (vert).
    Les quatre cellules affichent ici $C_L(\alpha)$, la polaire
    $C_L(C_D)$, la finesse $C_L/C_D$ en fonction de $C_L$ et le
    moment $C_M(C_L)$.}

\section{Configuration de la grille}

La grille est configurable via le menu
\menu{Affichage $\rightarrow$ Disposition}. Cinq dispositions
préconfigurées sont disponibles : $1\times1$, $1\times2$,
$2\times2$ (par défaut), $2\times3$, $3\times3$.

\begin{noteinfo}
Lors d'un changement de disposition, les analyses déjà
sélectionnées dans les cellules existantes sont conservées.
Les nouvelles cellules sont initialisées avec les analyses
par défaut.
\end{noteinfo}

\section{Cellule d'analyse}

Chaque cellule comporte une \textbf{barre de contrôle} et un
\textbf{canvas matplotlib}.

\subsection{Barre de contrôle}

\begin{tabularx}{\linewidth}{l X}
    \toprule
    \textbf{Élément} & \textbf{Rôle} \\
    \midrule
    Liste déroulante (analyse)
        & Choix du type d'analyse à afficher dans la cellule
          (voir \emph{Types d'analyse}, ci-après). \\
    Case \menu{C}
        & Affiche / masque les courbes du profil courant (bleu). \\
    Case \menu{R}
        & Affiche / masque les courbes du profil de référence (rouge). \\
    Case \menu{F}
        & Affiche / masque les courbes du profil avec volet (vert). \\
    Liste \menu{Reynolds}
        & Pour les analyses dépendant d'un Reynolds particulier
          ($-C_p(x)$, Cp~+~Profil, profil~+~CL), choisit le
          Reynolds à afficher. Masquée si non pertinente. \\
    Liste \menu{Alpha}
        & Pour les analyses dépendant d'une incidence (Cp~+~Profil,
          profil~+~CL et les grandeurs de couche limite), choisit
          l'incidence $\alpha$ à afficher. Masquée si non pertinente. \\
    Coordonnées
        & Affichage en temps réel des coordonnées du curseur
          dans le repère du graphique. \\
    \bottomrule
\end{tabularx}

\subsection{Interactions souris}

Les interactions disponibles sur le canvas matplotlib sont
identiques à celles de l'onglet \menu{Profils} :

\begin{itemize}
    \item Clic gauche + drag : zoom rectangle ;
    \item Molette : zoom centré sur le curseur ;
    \item Clic milieu : pan ;
    \item Clic gauche sur la légende : masque / réaffiche la
        courbe correspondante.
\end{itemize}

\subsection{Légende interactive et défilante}

Lorsqu'une cellule contient de nombreuses courbes (par exemple
plusieurs Reynolds combinés à plusieurs profils), la légende peut
devenir trop longue pour s'afficher entièrement. Une \textbf{barre
de défilement} apparaît automatiquement au-delà de 15~entrées.

\begin{astuce}
Pour comparer rapidement deux courbes, il est plus efficace
de \textbf{masquer} les courbes parasites en cliquant sur leur
entrée dans la légende, plutôt que de zoomer ou d'ajouter
une cellule supplémentaire.
\end{astuce}

\section{Types d'analyse disponibles}

\subsection{Polaires (en fonction d'Alpha)}

\begin{tabularx}{\linewidth}{l X}
    \toprule
    \textbf{Analyse} & \textbf{Description} \\
    \midrule
    $C_L(\alpha)$  & Coefficient de portance en fonction de
                     l'incidence. Une courbe par Reynolds.\\
    $C_D(\alpha)$  & Coefficient de traînée en fonction de
                     l'incidence.\\
    $C_M(\alpha)$  & Coefficient de moment de tangage
                     (référence quart de corde).\\
    $C_L/C_D$~$(\alpha)$ & \emph{Finesse} en fonction de
                     l'incidence. Permet de déterminer
                     l'incidence de finesse maximale.\\
    $C_L/C_D$~$(C_L)$ & \emph{Finesse} en fonction du coefficient de
                     portance. Pratique pour lire la finesse au
                     $C_L$ de croisière visé.\\
    $C_M(C_L)$     & Coefficient de moment en fonction de la
                     portance. La pente renseigne sur la stabilité ;
                     l'effet d'un volet y est très lisible.\\
    $C_D(C_L)$ \emph{(polaire de Lilienthal)}
                   & Polaire classique : traînée en abscisse
                     vs portance en ordonnée. La tangente
                     depuis l'origine donne la finesse maximale.\\
    Transition top, bot
                   & Position $x/c$ de la transition de couche
                     limite sur extrados / intrados, en fonction
                     de $\alpha$. Renseigne sur la rugosité
                     effective et le décrochage laminaire.\\
    \bottomrule
\end{tabularx}

\subsection{Distributions le long du profil}

\begin{tabularx}{\linewidth}{l X}
    \toprule
    \textbf{Analyse} & \textbf{Description} \\
    \midrule
    $-C_p(x)$ & Distribution du coefficient de pression le long
               de la corde, à un Reynolds et une incidence donnés.
               Convention : $-C_p$ vers le haut (extrados au-dessus). \\
    Cp + Profil & Tracé combiné façon \emph{XFoil} : le $-C_p$ en haut
               et le profil redessiné en dessous, pour un couple
               (Reynolds, incidence). Un encart résume les coefficients
               ($C_L$, $C_M$, $C_D$, finesse). Pour une vue dédiée plus
               complète (Cp non visqueux et couche limite), voir le
               chapitre~\ref{chap:cp}. \\
    Profil + CL & Visualisation combinée : profil tourné par
               l'incidence + tracé de la frontière de couche
               limite (épaisseur de déplacement $\delta^*$).\\
    \bottomrule
\end{tabularx}

\subsection{Couche limite}

Cinq grandeurs de couche limite peuvent être tracées en fonction
de l'abscisse curviligne $s$ ou de l'abscisse cartésienne $x$. Elles
sont évaluées à l'\textbf{incidence sélectionnée} (liste \menu{Alpha})
— la couche limite étant calculée pour chaque $\alpha$ — et chaque
courbe compare les Reynolds disponibles :

\begin{tabularx}{\linewidth}{l X}
    \toprule
    \textbf{Grandeur} & \textbf{Description} \\
    \midrule
    $U_e(s)$, $U_e(x)$
        & Vitesse au bord de la couche limite,
          normalisée par la vitesse infinie amont.\\
    $\delta^*(s)$, $\delta^*(x)$
        & Épaisseur de déplacement.\\
    $\theta(s)$, $\theta(x)$
        & Épaisseur de quantité de mouvement.\\
    $C_f(s)$, $C_f(x)$
        & Coefficient de frottement local.\\
    $H(s)$, $H(x)$
        & Facteur de forme $H = \delta^*/\theta$ ($H \approx 2{,}6$
          au point de transition, $H \approx 4$ au décollement).\\
    \bottomrule
\end{tabularx}

\section{Bandeau de résultats obsolètes}

Lorsque l'utilisateur revient à l'onglet \menu{Profils} et modifie
un profil après simulation (déplacement d'un point de contrôle,
changement de degré, conversion\dots), un bandeau jaune apparaît
en haut de l'onglet \menu{Résultats} :

\begin{quote}
\textcolor{orange!75!black}{$\triangle$
\emph{Le profil courant a été modifié depuis la dernière
simulation. Relancez pour mettre à jour.}}
\end{quote}

Ce bandeau disparaît automatiquement à la prochaine simulation.

\begin{important}
Les résultats affichés dans les cellules \emph{ne sont pas}
recalculés automatiquement après modification d'un profil.
Le bandeau est un simple indicateur visuel : il appartient
à l'utilisateur de relancer la simulation depuis l'onglet
\menu{Paramétrage XFoil}.
\end{important}
